% © 2026 Ing. David Jaroš — CC BY-NC-ND 4.0
%
% File: research_tracks/BRIDGE_CLOSURE/B4_rpsi_selfdual_vs_scale_bridge.tex
% Status: BRIDGE CLOSURE NOTE.
% Purpose: Correctly separate the proven self-dual R_psi=1 statement from the open physical scale calibration.

% UBT-AI-PROVENANCE-BEGIN
% schema: ubt-ai-provenance/v1
% tier: C_working
% ai_assistance: disclosed
% human_review: risk-based
% editorial_responsibility: Ing. David Jaroš
% policy: ../../AI_PROVENANCE.md
% notice: Working material; exhaustive human review is not claimed.
% UBT-AI-PROVENANCE-END

\documentclass[12pt,a4paper]{article}
\usepackage[a4paper,margin=2.5cm]{geometry}
\usepackage{amsmath,amssymb,booktabs}
\usepackage[most]{tcolorbox}
\usepackage{hyperref}
\usepackage{xcolor}
\usepackage{microtype}

\hypersetup{colorlinks=true,linkcolor=blue!60!black,urlcolor=blue!60!black,citecolor=green!50!black}

\newtcolorbox{statusbox}[1][]{
  colback=blue!5!white,
  colframe=blue!60!black,
  arc=3pt,
  boxrule=1pt,
  title=Bridge Status,
  #1
}

\title{Bridge B4: Self-Dual $R_\psi$ versus Physical Scale Calibration}
\author{Ing. David Jaroš}
\date{2026-06-14}

\begin{document}
\maketitle

\begin{abstract}
This note records a status correction.  The dimensionless self-dual radius
$R_\psi=1$ is not an open bridge: the repository already derives
$R_\psi=R_t$ from the self-dual/T-duality fixed point, hence $R_\psi=1$ in
natural units.  The remaining scale bridge is different: one must derive why
that self-dual unit corresponds to the physical unification scale
$M_{\rm UBT}\simeq2\times10^{16}$ GeV.
\end{abstract}

\begin{statusbox}
\textbf{Closed:}
\[
  R_\psi=R_t,
  \qquad R_t=1\Rightarrow R_\psi=1.
\]
\textbf{Open/MC:}
\[
  R_\psi=1
  \quad\Longrightarrow\quad
  M_{\rm UBT}\simeq2\times10^{16}\,\mathrm{GeV}.
\]
The open bridge is physical unit calibration, not self-dual radius fixation.
\end{statusbox}

\section{Closed part: self-dual radius}

The canonical geometry file
\texttt{canonical/geometry/Rpsi\_dynamical\_fix.tex} establishes the self-dual
condition
\[
  R_\psi = R_t.
\]
Choosing natural units for the real-time circle,
\[
  R_t=1,
\]
gives
\[
  \boxed{R_\psi=1}.
\]
This is the correct input for the alpha winding analysis and for the modular
self-dual point $\tau=i$.

\section{Open part: physical scale calibration}

The Weinberg RG branch needs a physical scale,
\[
  M_{\rm UBT}\simeq2\times10^{16}\,\mathrm{GeV}.
\]
This is not the same statement as $R_\psi=1$.  Rather, it requires a physical
unit map
\[
  M_{\rm UBT}=\Lambda_\psi/R_\psi.
\]
Since $R_\psi=1$, the actual bridge is
\[
  \Lambda_\psi\simeq2\times10^{16}\,\mathrm{GeV}.
\]
Thus the scale-closure task is to derive $\Lambda_\psi$, not to rederive
$R_\psi=1$.

\section{Current best candidate}

The file \texttt{research\_tracks/EW/rpsi\_from\_action.tex} records a viable
candidate based on Scherk--Schwarz half-mode projection and gauge-condensate
normalisation:
\[
  \lambda_{\rm eff}=\frac{g^2}{2}\approx0.122,
\]
while the numerical target is
\[
  \lambda_{\rm target}\approx0.119.
\]
The mismatch is about $2.7\%$, which is small enough to treat the route as a
serious candidate, but not yet an exact theorem.

\section{Corrected bridge wording}

The bridge-closure programme should no longer say:
\[
  \text{derive } R_\psi=1.
\]
It should say:
\[
  \text{derive the physical scale calibration }\Lambda_\psi
  \text{ at the self-dual point }R_\psi=1.
\]

\end{document}
